\documentclass{article} 
\input{../../lib.sty} 
\allowdisplaybreaks

\title{\texttt{fn function\_report\_top\_k}} 
\author{Michael Shoemate} 
\begin{document}  
\maketitle 
 
\section{Hoare Triple} 
\subsection*{Precondition} 
\subsubsection*{Compiler-verified}
Types consistent with pseudocode.

\subsubsection*{Caller-verified}
None
 
\subsection*{Pseudocode}
\lstinputlisting[language=Python,firstline=2,escapechar=|]{./pseudocode/function_report_top_k.py} 
 
\subsection*{Postcondition} 

\begin{theorem}
    Returns a noninteractive function with no side-effects that, when given a vector of scores $x_i$,
    returns the indices of the top k $y_i$,
    where each $y_i = -x_i$ if \texttt{optimize} is \texttt{min}, else $y_i = x_i$.

    The returned function will only error if there are no non-null scores.
\end{theorem}

\begin{proof}
    The comparator on line \ref{fn-optimize} effectively flips the sign of the scores if \texttt{optimize} is \texttt{min},
    therefore each $y_i = -x_i$ if \texttt{optimize} is \texttt{min}, else $y_i = x_i$.
    This is done to avoid the negation of a signed integer.

    Assume the scores are non-null, as required by the returned function precondition.
    Then \texttt{max\_sample} on line \ref{fn-max-sample} defines a total ordering on the scores.

    Since the score vector is finite, and \texttt{max\_sample} defines a total ordering,
    then the preconditions for \rustdoc{measurements/report_noisy_top_k/fn}{top} are met.
    Therefore on line \ref{fn-top} \texttt{top} returns the pairs with the top $k$ scores.
    Line \ref{fn-top} then discards the scores, returning only the indices,
    which is the desired output.

    There is one source of error in the function,
    when there are no non-null scores in the input vector.
\end{proof}
 
\end{document} 
